\subsection{Transformation and Informational Equivalence}

The generation of continuous spatial structures requires a specific
behavior from the Ground Reality. Based on the existence of \textbf{Causal
Assertions} ($c \in \mathbb{A}^\diamond$)---and the critical fact that they are
not immediately refuted by reality---the Subject operates under the assumption
of the \textit{Nice-ness} (regularity) of the Noumena. 

The Noumena tends to instantiate events that are `close' to existing History.
In informational terms, closeness implies that two points share massive
\textbf{Equivalence} and exhibit minimal \textbf{Differentiation} (e.g.,
differing by a magnitude of one minimal step in a Numeric Fact).

Because the Subject operates on pure information rather than a rigid physical
container a redundancy reducing \textbf{Transformation} (Projection) is
inherent. A synchronic set of points where \textit{all} $d$ discernibles vary
uniformly is informationally equivalent to a set of points where \textit{only
one} discernible varies. By nature, the Subject maps the complex co-variation into a
singular dimension of change.

\begin{formalism}[Orthogonality]
    \textbf{Orthogonality} is not a physical angle, but the \textit{Informational
    Independence} of discernibles within the composition of a Point. If a logical
    transformation can isolate the variation of one discernible such that its
    change does not necessitate a change in the others, the discernibles form
    an orthogonal basis. 
\end{formalism}

\subsection{Geometry as Transcendental Constraint}

With the equivalence of linear transformations established, the Subject generates
the conceptual space of Geometry. Geometric constructs are not arbitrary
inventions; they are the necessary fixed points of the Synthetic Interstice when
the generation of new Points ($p_{new}$) is restricted by the
\textbf{Transcendental Gate} (\texttt{WHERE} operator) against the truths of
History.

\begin{formalism}[Geometric Generation]
    A geometric object is the saturated set of points in the Interstice generated
    under a specific condition on History. 
    
    \textbf{The Line:} Generated when the \texttt{WHERE} constraint enforces a
    constant differential equivalence (slope) across the orthogonal discernibles.
    
    \textbf{The Circle:} Generated when the \texttt{WHERE} constraint enforces
    Equivalence on the \textit{Metric Fact} ($m$) between any newly generated
    point $p_{new}$ and a fixed historical anchor point $p_{anchor}$:
    
    \begin{equation}
        \text{Premise} \implies p_{new} \quad \textnormal{\texttt{WHERE}} \quad m(p_{new}, p_{anchor}) \equiv \text{constant}
    \end{equation}
    
\end{formalism}

Space, therefore, is not a pre-existing void. It is the \textbf{Implication of
Order} existing in the dynamic present, structured by the informational
transformations of the Subject, and constrained by the \texttt{WHERE} conditions
applied to the metric facts of History.

\section{The Genesis of Space}

For every discernible $\delta_x$---whether singular or composed---of a fact
$\langle \bot, \delta_x \rangle$ in History, there is a necessary antecedent:
it must have passed through the Synthetic Interstice carrying the truth of
\textit{Becoming} ($\top$). That is, at some point in the logical generation,
there must have existed an event $\langle \top, \delta_x \rangle$.

The events corresponding to sequential facts, namely $(\delta_1 \succ
\delta_0)$, establish the faculty of \textbf{Order} within the Synthetic
Interstice. Crucially, this order exists in the \textit{absence of time}, in a
gap of time (the Synthetic Interstice), to be precise. Thus, the ordering
itself \textit{implies} a new event: The Event of {\it Adjacency Bonding\/}
$e_{ab}$.

\begin{equation}
    e_{ab} = \langle \top, (\delta_1 \succ \delta_0) \rangle 
\end{equation}

We define the operator $\tilde{\implies}$ to express the operator of Order
insider the Synthetic Interstice.

It is \textit{identical} to the discernible of the sequential fact ($\succ$).
Space, therefore, is not a container, but the \textbf{Implication of Order}
existing in the dynamic present. We can now derive the {\it Adjacent Bonding
Fact\/} and the {\it Metric Fact}.

\begin{formalism}[Adjacent Bonding Fact]

    An \textbf{Adjacent Bonding Fact} $b \in \mathbb{A}^\top$ is the immutable
    assertion of `Close-ness' between two discernibles. It is the fixation of
    the immediate logical implication derived in the Synthetic Interstice.

    Let $\delta_1$ and $\delta_0$ be two discernibles linked by the proximity
    operator $\tilde{\implies}$. The Adjacent Bonding Fact is defined as:

    \begin{equation}
        b := \langle \bot, (\delta_1 \tilde{\implies} \delta_0) \rangle \in \mathbb{A}^\top
    \end{equation}

    It asserts that within the dynamic generation of the Now, the existence of
    $\delta_0$ was the immediate logical condition for $\delta_1$.

\end{formalism}

Just as the recursive composition of `After-ness' establishes the concept of a
Numeric Fact by means of the Sequential Fact, the recursive composition of
`Close-ness' establishes the concept of a \textbf{Metric Fact} by means of the
Adjacent Bonding Fact.

Again we consider two events that share the \textit{same} discernible $\delta$
within the same Synthetic Interstice, yet do not collapse into singularity. Let
them be distinguished by a chain of Adjacent Bonding Facts separating them.

\begin{formalism}[Metric Fact]

    A \textbf{Metric Fact} $m \in \mathbb{A}^\top$ is the immutable assertion
    of the magnitude of logical separation between two instances of an
    identical discernible $\delta$.

    It is established by the recursive composition of Adjacent Bonding Facts
    required to traverse from one instance to the other.

    \begin{equation}
        m := \langle \bot, (\tilde{\implies}^k, \delta) \rangle \in \mathbb{A}^\top
    \end{equation}

    Where $\tilde{\implies}^k$ represents the $k$-th degree of `Close-ness'.
    Metric Facts constitute the order coordinate of an entity relative
    to its identical counterpart.

\end{formalism}

The \textbf{Metric Fact} is the spatial counterpart to the \textbf{Numeric
Fact}. While the Numeric Fact anchors the identity of an entity in
\textit{Time}, the Metric Fact anchors the identity of an entity in a dimension
`orthogonal' to time.\ Let this dimension be called  \textit{Space}. 

-------------
The $\varnothing$ discernible is the primal experience of Non-Truth. When
ignored, it carries the essence of \textit{Illusion}: the belief that a
`causal reason'---which is a concept---exists ontologically, while all
conceptual experience actually emerges through the Synthetic Interstice
triggered by events instantiated from the inconceivable Noumena. When the
$\varnothing$ discernible is ignored, the Cognitive Subject enters the
illusory world of pure conceptualization.


Instantiations of events have their cause in the Noumena. Implications are
fixated rules. To exist, they must have been instantiated as events, which then
where fixated into History. Since, the history is immutable, those rules
themselves are immutable.

\section{Adjacency and Metric Facts}

We established the \textbf{Numeric Fact} as the measure of the recurrence of
Identity across the boundary of the Fixation ($\top$ vs $\bot$). We now derive
the \textbf{Metric Fact} as the measure of the separation of Identity across
the boundary of logical Necessity.

In the Synthetic Interstice, two events $e_a$ and $e_b$ with a shared
discernible cannot coexist except if they posses at least one different
discernible. Otherwise, they are equivalent and collapse into a single
identity. For being different in on discernible, they need to be composed,
i.e.\ at least one of them is necessarily the result of an {\it Implication}. 

\begin{formalism}[Adjacency Fact]

    An \textbf{Adjacency Fact} $j \in \mathbb{A}^\top$ is the immutable record
    of a direct implication between two discernibles within a single Synthetic
    Interstice.
    
    If, during the genesis of $N_k$, the existence of $\delta_a$ implied
    $\delta_b$ (i.e., $\delta_a \implies \delta_b \in \mathcal{I}$), this
    relation is fixated as:

    \begin{equation}
        j := \langle \bot, (\delta_b \perp \delta_a) \rangle
    \end{equation}

    where $\perp$ denotes \textbf{Structural Adjacency}.

\end{formalism}

A third event $e_c$ needs to distinguish from $e_a$ and $e_b$ at least by one
discernible, so it requires another level of implication. The double
implication becomes its discernible. Just as the Sequential Fact ($\succ$)
orders History, the Adjacency Fact ($\perp$) orders the internal structure of
the Now.

\begin{formalism}[Metric Fact]

    A \textbf{Metric Fact} $m \in \mathbb{A}^\top$ is a record of applied
    cardinality of Implication. It establishes the magnitude of separation
    between two equal discernibles $\delta$ in the same Fixed Point.
    
    It is defined recursively as the count of Adjacency steps required to
    traverse from one instance of $\delta$ to the other:

    \begin{equation}
        m_{ab} := \langle \bot, (\perp^k, \delta) \rangle \in \mathbb{A}^\top
    \end{equation}
    
    Here, $\perp^k$ represents the $k$-th degree of logical separation
    (Implication Depth) in the Synthetic Interstice.

\end{formalism}

This provides a purely logical derivation of Time and Space:

\begin{itemize}
    \item \textbf{Time} is the cardinality of Fixations (Vertical Depth).
    \item \textbf{Space} is the cardinality of Implications (Horizontal Depth).

\end{itemize}

\textbf{Metric} is therefore not a physical property of a container, but the
\textbf{Latency of Logic}. Two points are ``far apart'' simply because the
Manifold requires many logical steps to derive one from the other.

\section{Topology}

TODO: self by consideration of not self concept + Fact of Belonging (by one discernible) + Metric Fact

Since the faculty of Equivalence and Differentiation is already established in
the Synthetic Interstice, Implications can establish events that represent
`adjacency' between events like succession does in the sequence of History.
Topology can now be established as a `causality' orthogonal to primitive time.

\begin{formalism}[Topological Event]

    A \textbf{Topological Event} $e_t \in \mathbb{E}$ is an ephemeral entity
    emerging within the Synthetic Interstice that manifests the mutual
    implication between two co-existing events.

    Let $\bowtie$ denote the adjacency operator (representing structural
    ``closeness''). If two distinct events $e_1 = \langle \top, \delta_1
    \rangle$ and $e_2 = \langle \top, \delta_2 \rangle$ imply each other via
    the Implication Set $\mathcal{I}$ during the interstice:

    \begin{equation}
        (e_1 \implies e_2) \quad \land \quad (e_2 \implies e_1)
    \end{equation}

    then an implication event $e_t$ is generated:

    \begin{equation}
        e_t := \langle \top, (\delta_1 \bowtie \delta_2) \rangle
    \end{equation}

    The event $e_t$ exists purely in the domain of Becoming ($\top$) and
    carries the composite discernible $(\delta_1 \bowtie \delta_2)$. It serves as
    the fundamental unit of dynamic topological extension prior to any historical fixation.

\end{formalism}

\begin{formalism}[Topological Assertion]

    A \textbf{Topological Assertion} $t \in \mathbb{A}^\diamond$ is the speculative fixation
    that, behind a recurring network of Topological Facts $p \in \mathbb{A}^\top$, there
    exists a non-conceptual reality in the Noumena that governs this synchronous regularity.

    Let $p = (\delta_1 \bowtie \delta_2)$ be a Topological Fact. The Topological Assertion
    $t \in \mathbb{A}^\diamond$ posits that this adjacency is an intrinsic rule of the
    metaphysical substance, defined as:

    \begin{equation}
        t := (p \xrightarrow{\mathcal{N}} \top)
    \end{equation}

    Like the Causal Assertion, it is refutable; it remains valid only as long as the
    Synthetic Interstice continues to reliably produce the adjacency $p$ in future states.

\end{formalism}
